\section{Enforcing Non-Negativity}\label{sec:primaldual}

Section~\ref{sec:pd} reduces the long-only problem to solving an SPD system, leaving
$w \geq 0$ as the sole remaining constraint.  Two strategies are available.

\paragraph{Active-set strategy}
The first strategy separates the budget constraint from the non-negativity constraint.
The budget-only problem is solved over a dynamically adjusted \emph{active set} of
assets; non-negativity is then enforced iteratively without ever modifying the inner
solve.

In the \emph{primal step}, any asset with a negative weight is dropped from the active
set and the budget-constrained SPD system is re-solved on the reduced universe.  Once
all active weights are non-negative, the \emph{dual step} checks the KKT gradient
condition for every excluded asset: global optimality requires every excluded asset $i$
to satisfy
\[
  \nabla_i f(w) \;\geq\; \lambda\;,
  \qquad
  \nabla_i f(w) = \tfrac{2}{T}(X^\top Xw)_i - \rho\hat\mu_i\;,
\]
where $\lambda$ is the budget Lagrange multiplier recovered analytically from
Section~\ref{sec:pd}.  An excluded asset that violates this condition would decrease
portfolio variance if held at a small positive weight; it is re-added to the active set
and the loop repeats.  The loop terminates when both primal and dual feasibility hold;
together with stationarity from the inner solve this certifies global optimality.

\begin{algorithm}[ht]
\caption{Active-set outer loop (wraps any inner solver $f$)}\label{alg:elim}
\begin{algorithmic}[1]
\Require Inner solver $f$;\; tolerance $\varepsilon = 10^{-6}$
\Ensure Weight vector $w \in \mathbb{R}^n$;\; total inner iteration count $k$
\State $\mathrm{active} \leftarrow \{1,\ldots,n\}$;\; $k \leftarrow 0$
\Loop
  \State $(w_{\mathrm{a}},\, k_{\mathrm{step}}) \leftarrow f(\mathrm{active})$;\; $k \leftarrow k + k_{\mathrm{step}}$
  \If{$\exists\, i \in \mathrm{active}$ with $w_{a,i} < -\varepsilon$}
    \State $\mathrm{active} \leftarrow \mathrm{active} \setminus \{i : w_{a,i} < -\varepsilon\}$
      \hfill\textit{(primal step: drop violators)}
  \Else
    \State $w_i \leftarrow w_{a,i}$ for $i \in \mathrm{active}$;\; $w_i \leftarrow 0$ for $i \notin \mathrm{active}$
    \State $g_i \leftarrow \tfrac{2}{T}(X^\top Xw)_i - \rho\hat\mu_i$
    \State $\lambda \leftarrow \mathrm{mean}_{j \in \mathrm{active}}(g_j)$
    \State $\mathcal{V} \leftarrow \{i \notin \mathrm{active} : g_i < \lambda - \varepsilon\}$
    \If{$\mathcal{V} = \emptyset$}
      \State \textbf{break} \hfill\textit{(KKT satisfied; globally optimal)}
    \EndIf
    \State $\mathrm{active} \leftarrow \mathrm{active} \cup \mathcal{V}$
      \hfill\textit{(dual step: re-add violators)}
  \EndIf
\EndLoop
\State \Return $w$, $k$
\end{algorithmic}
\end{algorithm}

\paragraph{Simplex-projection strategy}
The second strategy treats the two constraints as a single geometric object: the
probability simplex
$\Delta_n = \{w \in \mathbb{R}^n : \mathbf{1}^\top w = 1,\, w \geq 0\}$.
Rather than separating budget and non-negativity, gradient steps of the form
\[
  w_{k+1} = \Pi_{\Delta_n}\!\left(w_k - \tau\,\nabla f(w_k)\right)
\]
project directly onto $\Delta_n$ after each descent step, where $\Pi_{\Delta_n}$
is the Euclidean projection.
Duchi et al.\ \cite{duchi2008} and Condat~\cite{condat2016} show that $\Pi_{\Delta_n}$ can be computed in
$O(n\log n)$ by a single sort-and-threshold pass, making each iteration cheap.
There is no outer loop: both constraints are handled simultaneously, and the method
terminates when the step size $\|w_{k+1} - w_k\|$ falls below a tolerance.

The gradient $\nabla f(w) = \tilde{X}^\top(\tilde{X}w) - \tilde{X}^\top\mathbf{0}
= \tilde{X}^\top(\tilde{X}w)$ is evaluated matrix-free in $O(nT)$ using two
matrix-vector products with $\tilde{X}$, at the same cost per step as the CG inner
solver.  The Lipschitz constant $L = \lambda_{\max}(\tilde{X}^\top\tilde{X})$, needed
for the step size $\tau = 1/(2L)$, is estimated once by power iteration on $\tilde{X}$,
again matrix-free.

\begin{algorithm}[ht]
\caption{Simplex-projection solver (proximal gradient, matrix-free)}\label{alg:proximal}
\begin{algorithmic}[1]
\Require Stacked matrix $\tilde{X} \in \mathbb{R}^{(T+n_R)\times n}$;\;
         tolerance $\varepsilon = 10^{-6}$;\; $k_{\max} = 1{,}000$
\Ensure Weight vector $w \in \Delta_n$
\State Estimate $L \leftarrow \lambda_{\max}(\tilde{X}^\top\tilde{X})$ by power iteration
  \hfill\textit{(matrix-free, $O(n_{\mathrm{pow}}\cdot nT)$)}
\State $\tau \leftarrow (2L)^{-1}$;\; initialise $w_0 \in \mathbb{R}^n$ randomly;\; $k \leftarrow 0$
\Loop
  \State $g \leftarrow \tilde{X}^\top(\tilde{X}w_k)$
    \hfill\textit{(matrix-free gradient, $O(nT)$)}
  \State $w_{k+1} \leftarrow \Pi_{\Delta_n}(w_k - \tau\,g)$
    \hfill\textit{(simplex projection, $O(n\log n)$)}
  \If{$\|w_{k+1} - w_k\|_2 \leq \varepsilon$ \textbf{or} $k \geq k_{\max}$}
    \textbf{ break}
  \EndIf
  \State $k \leftarrow k + 1$
\EndLoop
\State \Return $w_{k+1}$
\end{algorithmic}
\end{algorithm}

\medskip
The two strategies make different tradeoffs. The active-set loop solves a sequence of
linear systems and benefits from fast inner solvers (Section~\ref{sec:solvers}); its
convergence to global optimality is certified by the KKT check. The simplex-projection
approach is single-stage and loop-free, but converges in $O(\kappa)$ gradient steps
versus $O(\sqrt{\kappa})$ for CG, so it is most competitive when the condition number is
moderate or the matrix-free $O(nT)$ cost dominates.
Section~\ref{sec:solvers} develops concrete implementations of both strategies.
